\subsection{一个由对应构成的范畴}\label{SS: cat of corr}
设 $(\overleftarrow{c},\overrightarrow{c}): C\rightrightarrows X$ 是空间中的群胚。特别地，有下图
\begin{equation*}
\label{E: diagram}
\xymatrix{
&& C\times_XC\ar_{\overleftarrow{p}}[ld]\ar^{\overrightarrow{p}}[dr]\ar^m[d]& &\\
&C\ar_{\overleftarrow{c}}[ld]\ar_{\overrightarrow{c}}[rd]& C\ar_{\overleftarrow{c}}[lld]\ar^{\overrightarrow{c}}[rrd]&C\ar^{\overleftarrow{c}}[ld]\ar^{\overrightarrow{c}}[rd]&\\
X&&X&&X
}
\end{equation*}
其中 $\overleftarrow{p},\overrightarrow{p}: C\times_XC\to C$ 是第一、第二因子的投影，$m: C\times_XC\to C$ 是乘法映射。另外有映射 $\delta: X\to C$，使 $(\overleftarrow{c},\overrightarrow{c})\circ \delta:X\to X\times X$ 是对角映射。

现假设 $m$ 紧合且可由代数空间表示，$\delta$ 是闭浸入。
如下定义 $E$-线性范畴 $\on{D}^C(X)$：对象是 $X$ 上的层；态射为
\begin{equation}
\label{E:HomDC}
\Hom_{\on{D}^C(X)}(\mF_1,\mF_2)=\Hom_{\on{D}(C)}(\overleftarrow{c}^*\mF_1, \overrightarrow{c}^!\mF_2),
\end{equation}
即支撑在 $C$ 上、从 $\mF_1$ 到 $\mF_2$ 的上同调对应。

上一小节回顾的上同调对应形式化允许定义态射的合成律
\begin{eqnarray*}
&&\Hom_{\on{D}^C(X)}(\mF_1,\mF_2)\otimes\Hom_{\on{D}^C(X)}(\mF_2,\mF_3)\\
&=&\on{Corr}_C((X,\mF_1),(X,\mF_2))\otimes \on{Corr}_C((X,\mF_2),(X,\mF_3))\\
&\to& \on{Corr}_{C\times_XC}((X,\mF_1),(X,\mF_3)) \\
&\xrightarrow{m_!}& \on{Corr}_C((X,\mF_1),(X,\mF_3))\\
&=&\Hom_{\on{D}^C(X)}(\mF_1,\mF_3).
\end{eqnarray*}
它含一个全子范畴 $\on{P}^C(X)$，由那些在 $X$ 上反常的 $\mF\in\on{D}^C(X)$ 组成。

考虑上述范畴的动机如下。
设 $f:X\to Y$ 是可表示的紧合态射，$C=X\times_YX$，$\overleftarrow{c}$ 与 $\overrightarrow{c}$ 是两投影。
按定义，
\begin{lem}
\label{L: fully faithful}
函子 $f_!: \on{D}(X)\to\on{D}(Y)$ 分解为
\[\on{D}(X)\to \on{D}^C(X)\to \on{D}(Y)\]
且第二个箭头全忠实。
\end{lem}
\begin{proof}
第一个断言是显然的。
第二个断言来自典范同构（由紧合基变换）
\[\Hom(\overleftarrow{c}^*\mF_1,\overrightarrow{c}^!\mF_2)=\Hom(f_!\mF_1,f_!\mF_2),\]
它与合成相容。
\end{proof}

\begin{rmk}
\label{R: DC enh}
(1) 作上述定义的一个理由在于：有些情形下即使叠 $Y$（从而 $\on{D}(Y)$）性态很差，范畴 $\on{D}^C(X)$ 仍可有恰当的定义。

(2) 上面作为 $E$-线性范畴引入的 $\on{D}^C(X)$ 已足够本文的应用。不过我们说明如何把它定义为 $\infty$-范畴，至少在 $\overleftarrow{c}$ 与 $\overrightarrow{c}$ 都是有限或代数闭域 $k$ 上分离（完全）有限型概形的（完全）紧合态射的情形。记 $\Shv(-)$ 为 $\on{D}(-)$ 的自然 $\infty$-范畴加强。若 $X$ 在 $\ell$-上同调维数有限的域上分离且（完全）有限型，则当 $E=\bF_\ell,\bZ_\ell, \bQ_\ell$ 等时，$\Shv(X,E)$ 由紧对象生成，紧对象即可构造复形。

注意 $\on{D}^C(X)$ 是 $E$-线性范畴 $2$-范畴中下图之余极限的单子式描述
\[\on{D}(C\times_XC)\ \substack{\longrightarrow\\[-1em] \longrightarrow \\[-1em] \longrightarrow}\ \on{D}(C)\ \substack{\longrightarrow\\[-1em] \longrightarrow} \ \on{D}(X),\]
其中函子由紧合前推给出。
于是显然可以把 $\infty$-范畴 $\Shv^C(X)$ 定义为
\[\cdots \substack{\longrightarrow\\[-1em] \longrightarrow \\[-1em] \longrightarrow\\[-1em] \longrightarrow}\ \Shv(C\times_XC)\ \substack{\longrightarrow\\[-1em] \longrightarrow \\[-1em] \longrightarrow}\ \Shv(C)\ \substack{\longrightarrow\\[-1em] \longrightarrow} \ \Shv(X)\]
关于紧合前推的余极限。
由 Barr–Beck–Lurie 定理，$\Shv^C(X)$ 有关于 $\Shv(X)$ 的单子式描述。特别地，其同伦范畴的 Hom 空间可用上同调对应描述：设 $\mF_1,\mF_2\in \Shv(X)$，经自然函子 $\Shv(X)\to \Shv^C(X)$ 看作 $\Shv^C(X)$ 的对象。则
$$\Hom_{\on{h}\Shv^C(X)}(\mF_1,\mF_2)=R\Hom_{\on{D}(C)}(\overleftarrow{c}^*\mF_1,\overrightarrow{c}^!\mF_2).$$
若 $C=X\times_YX$，$f:X\to Y$ 是某紧合满射，则 $f_!:\Shv(X)\to \Shv(Y)$ 诱导 $\infty$-范畴的等价 $\Shv^C(X)\simeq \Shv(Y)$。
\end{rmk}
